% THETA3-B21 contribution by /root/definite_bridge, 21 November 2026.
\subsection{THETA3-B21: the cross-type scalar is a specific determinant period}\label{B21}

B19 gives a common scalar when the complete CM types are fixed.
For different complete types, the remaining comparison can be
specified more precisely.  Balancing the four original CM types
controls the selected tangent periods, but need not control the
ordinary determinant of the full Hodge space of each theta torus.
This distinction does not disprove an equality between particular
period numbers.  It identifies the additional comparison that
the proposed CM-Hodge argument must actually establish.

\paragraph{An unambiguous modular normalization.}
Let $A_i/\overline\Q$ be the actual principally polarized CM
torus of B19, after its recorded isogeny, and set
$g=3[F:\Q]$.  Fix an algebraic basis of regular differentials
and normalized analytic coordinates by
\[
 A_i(\C)=\C^g/(\Z^g+Z_i\Z^g),\qquad
 \boldsymbol\omega_i=\Omega_i\,d\boldsymbol w,
 \qquad D_i=\det\Omega_i\ne0.
 \tag{B21a}
\]
Thus $\Omega_i$ is a period block of algebraic differentials,
not the normalized Siegel matrix $Z_i$.
Write $a\sim_{\overline\Q}b$ for
$a/b\in\overline\Q^\times$.
Use the geometric modular-form convention
\[
 f(A,\omega)=\frac{f(A)}{\omega^{\otimes k}},\qquad
 f(A,\lambda\omega)=\lambda^{-k}f(A,\omega).
\]
The analytic form $f_{\rm an}$ corresponding to the geometric
weight-$k$ form satisfies
\[
 f(A_i,\omega_{i1}\wedge\cdots\wedge\omega_{ig})
       =\frac{(2\pi i)^{gk}}{D_i^k}f_{\rm an}(Z_i).
 \tag{B21b}
\]
This is Lachaud--Ritzenthaler--Zykin,
\emph{Jacobians among abelian threefolds},
\S2.1, equation (1), Proposition 2.2.1 and Proposition 2.4.4
with its proof, printed pp.3--6, arXiv:0802.4017v1.
Their normalized block is called $\Omega_2$.
Although their application is in genus three and their displayed
spaces have full level, the proof of (B21b) is a change of
coordinates on a polarized complex torus and is unchanged after
fixing a finite level structure.

Here is the precise theta input.  Kirschmer--Narbonne--Ritzenthaler--Robert,
\emph{Spanning the isogeny class of a power of an elliptic curve},
arXiv:2004.08315v2, \S4.3, equation (7), Definition 4.3 and
Remark 4.4, printed pp.23--24, prove that pair products of the
finite-level theta-null coordinates are algebraic weight-one
Hodge sections.  They obtain this from the theta transformation
law and the algebraic $q$-expansion principle.  A basis of
algebraic differentials therefore fixes their modular affine
lift up to a common sign.  We use the scalar convention
(B21b) to fix the direction of the period comparison.

The theta coordinates in that source and the coordinates of B19
differ by a finite Fourier transform with algebraic entries.
Indeed, putting $n=N_i^2$ and using $a,j\in(\Z/n\Z)^g$,
\[
 \sum_j e^{-2\pi i a^tj/n}
     \vartheta\left[\begin{smallmatrix}0\\j/n\end{smallmatrix}\right]
                       (w,Z_i/n)
   =n^g\vartheta\left[\begin{smallmatrix}a/n\\0\end{smallmatrix}\right]
                       (nw,nZ_i).
 \tag{B21c}
\]
To prove this identity, expand the left side over $m\in\Z^g$.
The inner character sum is $n^g$ when $m\equiv a\pmod n$
and zero otherwise.  Writing $m=n\ell+a$ gives the right side.
At $w=0$ the completion factor in (B19b) is one.
Thus the pair products of B19's theta-null values have the
same algebraic modular interpretation.  Enlarging the finite
level to kill a multiplier, or taking a common even power,
would give the same conclusion below after taking an algebraic
root; no claim that the eighth power suffices for every chosen
rational characteristic is needed.

This interpretation is on the actual $A_i$, with its recorded
theta level.  In the coordinates of (B21c), the torus with
matrix $Z_i/n$ is the quotient of $A_i$ by the indicated
finite subgroup of its $n$-torsion; the map is induced by
the identity on $\C^g$.  Its algebraic model and this isogeny
are defined over $\overline\Q$, since the subgroup consists
of algebraic torsion points.  Pullback identifies their Hodge
spaces over $\overline\Q$.  Thus one may either use the
level-$n$ Hodge section on $A_i$ directly, as in the cited
theta-moduli construction, or pull it back through this
isogeny.  The latter description changes an algebraic
differential basis by a matrix in
$\operatorname{GL}_g(\overline\Q)$, whose determinant is
the recorded algebraic isogeny factor.  It introduces no
additional transcendental scalar and does not replace $D_i$
by an unspecified period of another abelian variety.

Let $\upsilon_i$ denote any scalar by which the specified
oscillator normalization multiplies the displayed standard
theta sum of (B19b); for exactly that displayed sum
$\upsilon_i=1$.  Define $\widetilde c_i=c_i/\upsilon_i$.
Then
\[
             \widetilde c_i^{\,2}
                    \sim_{\overline\Q}
                         \frac{D_i}{(2\pi i)^g}.
 \tag{B21d}
\]
To prove this, select an algebraic finite input $\varphi$ for
which $J_i\varphi(0)\ne0$.  B19 gives such an input and
identifies $\vartheta_\varphi(0)
=\widetilde c_i J_i\varphi(0)$ in the fixed rigidification.
The square of this theta-null is the value of an algebraic
weight-one modular form in analytic coordinates.  By (B21b),
\[
 \frac{(2\pi i)^g}{D_i}
       \vartheta_\varphi(Z_i,0)^2\in\overline\Q^\times.
\]
Dividing by the nonzero algebraic number
$J_i\varphi(0)^2$ proves (B21d).
If a higher even power was used, the assertion follows because
$\overline\Q$ is algebraically closed.  This proof retains
both the period orientation and the universal power of $2\pi i$.

\paragraph{The exact residual factor for two balanced pairs.}
For pairs $(1,2)$ and $(3,4)$, write B19's constants as
\[
 C_{x,ij}=\kappa_{x,ij}
                    \frac{c_ic_jp_ip_j}{N_iN_j}.
\]
Define
\[
 \mathcal R=
 \frac{C_{x,12}}{C_{x,34}}
 \frac{\kappa_{x,34}}{\kappa_{x,12}}
 \frac{\upsilon_3\upsilon_4}{\upsilon_1\upsilon_2}
 \frac{p_3p_4}{p_1p_2}\frac{N_1N_2}{N_3N_4}.
\]
Then $\mathcal R=\widetilde c_1\widetilde c_2/
(\widetilde c_3\widetilde c_4)$ exactly, and
\[
              \mathcal R^2\sim_{\overline\Q}
                        \frac{D_1D_2}{D_3D_4}.
 \tag{B21e}
\]
The universal powers cancel as
$(2\pi i)^{-g-g+g+g}=1$ because all four theta tori have
the same dimension.  The level factors have been displayed
explicitly; they are nonzero rational numbers.

Suppose $d_1+d_2=d_3+d_4$ as complete type functions,
with the indexing $d_i(\tau)=v_i(\tau^{-1})$ of B19.
The selected tangent eigenspace has a Hodge type which is
linear in $d_i$: at a conjugate embedding it is determined
by $(1+d_i(\tau))/2$, or by its simultaneous dual convention.
Hence the two selected tangent tensor products are isomorphic
as coefficient-field Hodge structures.  The accepted CM Hodge
conjecture makes this tensor comparison algebraic: after the
usual K\"unneth projectors it is a Hodge class on a product of
CM abelian varieties, represented by an algebraic cycle.
Its Betti and de Rham matrices are algebraic, so compatibility
of the comparison isomorphism gives
\[
                         \frac{p_1p_2}{p_3p_4}
                                      \in\overline\Q^\times.
 \tag{B21f}
\]
The factors $p_i$ here have the exact direction fixed by
(B19d); taking all four duals changes neither this conclusion
nor the balancing condition.

There is no hidden inverse $D_i$ in the Betti constant-line
calculation of 032, (4.6).  At the rational point $x=[Ez]$,
the line $L$, its annihilator, and every full compact embedding
summand in (4.5) are defined over $\overline\Q$ in the
complexification of the rational symplectic space.  A rational
symplectic change of coordinates is algebraic, and the graph
matrix $Z_i$ and its graph basis have algebraic entries.
The projection defining this graph is invertible by positivity.
The isomorphism
$\det Q^*=L\otimes\det(V_\mu)^*$ and the tensor determinant
identities use algebraic matrices.  Algebraic square roots can
be chosen in $\overline\Q$.  The ball coordinates obtained
from $E(z+z_1e_1+z_2e_2)$ give an algebraic Betti tangent frame
as well.  Thus choosing these algebraic constant-line frames
makes $\kappa_{x,12}/\kappa_{x,34}$ algebraic after the same
rank- and Taylor-order-dependent oscillator constants are fixed.
The original 032, p.15 immediately after (4.6), permits arbitrary
nonzero complex constant-line vectors and explicitly asserts no
rational trivialization; any such additional rescalings must
instead remain in the displayed $\kappa$ ratio.

Consequently, under the stated algebraic Betti frames, common
oscillator convention, and accepted CM-Hodge comparison,
\[
 \left(\frac{C_{x,12}}{C_{x,34}}\right)^2
       \sim_{\overline\Q}\frac{D_1D_2}{D_3D_4},
 \qquad
 \frac{C_{x,12}}{C_{x,34}}\in\overline\Q^\times
 \Longleftrightarrow
 \frac{D_1D_2}{D_3D_4}\in\overline\Q^\times.
 \tag{B21g}
\]
This is a reduction to one specified period ratio, not an
assertion that the ratio is algebraic.

\paragraph{Why complete-type balance does not settle that ratio.}
Assume for this calculation that $E/\Q$ is Galois, as in the
type-function setting of 032.  Put $h_d=(1+d)/2$ on
$\operatorname{Gal}(E/\Q)$ and let
\[
 A_d(\gamma)=\sum_\tau h_d(\tau)h_d(\gamma\tau).
 \tag{B21h}
\]
The ordinary top Hodge determinant of $B(d)$ is formed by
wedging the eigenlines indexed by $h_d(\tau)=1$.  Its
coefficient-field Betti determinant line has, at the conjugate
embedding $\gamma$, holomorphic degree $A_d(\gamma)$:
the sum simply counts which of the selected eigenlines remain
of type $(1,0)$.  By the actual decomposition (B19a), the
corresponding degree function for $A_{x,i}$ is
\[
                         M_d=2A_d+A_{d^{\rm flip}}.
 \tag{B21i}
\]
This is a quadratic function of the original type, unlike the
selected tangent type in (B21f).

The determinant $D_i$, up to an algebraic factor, is the
period of this selected determinant tensor.  Indeed choose
algebraic CM eigen-differentials.  Each row of their period
matrix is its one CM comparison period times a row with
algebraic entries in a rational Betti basis.  Taking a
nonzero period-block determinant gives the product of those
comparison periods times a nonzero algebraic determinant.
Algebraic changes of differential basis and the recorded
isogenies preserve this statement.

For a direct counterexample to the proposed inference, take
$E=\Q(\zeta_7)$ and order its cyclic Galois group as
$1,\gamma,\ldots,\gamma^5$, with complex conjugation
$\gamma^3$.  Let $\mu=1$ and define the four complete types
\[
 d_{a,b}=(1,a,b,-1,-a,-b),\qquad a,b\in\{1,-1\}.
\]
They satisfy $d_{+,+}+d_{-,-}=d_{+,-}+d_{-,+}$, and all
have the same positive value at $\mu$.
Flipping $\mu,c\mu$ changes entries zero and three.
Direct counting in (B21h) gives
\[
\begin{array}{c|rrrrrr}
 d&M_d(1)&M_d(\gamma)&M_d(\gamma^2)&M_d(\gamma^3)
       &M_d(\gamma^4)&M_d(\gamma^5)\\\hline
 d_{+,+}&9&6&3&0&3&6\\
 d_{-,-}&9&6&3&0&3&6\\
 d_{+,-}&9&4&5&0&5&4\\
 d_{-,+}&9&2&7&0&7&2
\end{array}
\]
For example $A_{+,+}=(3,2,1,0,1,2)$, whereas
$A_{-,+}=(3,0,3,0,3,0)$; flipping interchanges the relevant
entries as prescribed.  Subtraction yields
\[
 M_{+,+}+M_{-,-}-M_{+,-}-M_{-,+}
                         =(0,6,-6,0,-6,6)\ne0.
 \tag{B21j}
\]
Thus the two specified determinant tensors need not be
isomorphic as coefficient-field Hodge structures, even though
the original complete types balance.  Applying the CM Hodge
conjecture to a nonexistent Hodge isomorphism cannot prove
their period ratio algebraic.

If the stronger equality $M_{d_1}+M_{d_2}=M_{d_3}+M_{d_4}$
does hold, the same algebraic-cycle argument used for the
tangent tensors proves the determinant period ratio algebraic.
Without it, (B21j) invalidates that particular deduction; it
does not prove the period ratio transcendental, exclude other
period relations, or refute the original mixed-pairing theorem.
The exact finite-jet functional and the common normalization
for fixed complete types in B19 remain valid.\hfill$\square$
